\section{Assembly at a fixed curve}\label{pw:section}

Fix $E/\mathbb Q$ with nonzero rational two-torsion.  The even construction
has two uses: it proves the corank-zero case and supplies the uniform
estimate needed to normalize the odd coefficients.  After that normalization,
the corank-one argument fixes a nonvanishing companion and constructs the
two simultaneous coefficient tests required for ring-class transfer.

\subsection{Interfaces and conventions}\label{pw:interfaces-section}

We use the fixed-support conventions of Section~\ref{sec:preliminaries}.
Thus all variable prime factors are odd primes of good reduction,
$p^*=(-1/p)p$, and a filter fixes the real sign and unit squareclasses
at a finite support containing $2N$, with finitely many further exclusions.
Recall
\[
 t_p=\dim_{\mathbb F_2}E[2]^{\operatorname{Fr}_p},\qquad
 w(h)=\frac12\sum_{p\mid h}t_p.
\]
Whenever the two-primary group is finite, we write
\[
s(h)=\log_2|\Sha(E^{(h)}/\mathbb Q)[2^\infty]|.
\]
This includes a base of Selmer corank zero and every parameter of analytic
order zero or one.  It is distinct from $s_2(E^{(h)})$, the Selmer corank.
The normalized central value is denoted $L_*(h)$, and $v=v_2$ with
$v(2)=1$.

Here are the preceding results in the form used below.  Their roles are
different: coefficient units give upper bounds, interpolation returns those
bounds to an unknown base, and the ring-class lower bound makes the odd
normalization possible.  The latter step is carried out in
Proposition~\ref{pw:odd-lower}.

\begin{proposition}[Collected interfaces]\label{pw:interfaces}
The following statements hold for the fixed curve and fixed local data.
\begin{enumerate}
\item The cyclotomic missing-vertex statement of
Theorem~\ref{cy:missing-vertex} applies to any fixed base $h_0$ of Selmer
corank zero.  If arbitrarily large binary cubes based at $h_0$, with no old
prime removed, have nonzero central values at every other address and satisfy
\[
 v(L_*(h_x))-2w(h_x)-s(h_x)\le B\quad(x\ne0),
\]
then the base central value is nonzero and satisfies the same inequality
with $B$ replaced by $B+C$.  The constant $C$ depends only on the fixed curve
and local data, not on $h_0$, its number of prime factors, or the cube
dimension.  The dimension required for the argument may depend on the base.
The cyclotomic lower bound at an analytically nonvanishing parameter is also
available.
\item Corollary~\ref{rd:ring-lower} has the following uniform form.  If
$k<0$ is an odd fundamental discriminant with at most $A$ prime factors,
$K=\mathbb Q(\sqrt{k})$ is neither an extra-unit field nor the CM field,
$k$ is a square at $2Nh$, and the twists $h,hk$ have respective analytic
orders one and zero, then their ring-class Heegner sum has free-line index
$j_h(k)$ satisfying
\begin{equation}\label{pw:ring-lower-interface}
 2j_h(k)\ge 4w(h)+s(h)+s(hk)-C_{E,A}.
\end{equation}
For a fixed such $k$, Theorem~\ref{rd:missing-vertex} applies to
cubes based at $h=1$ whose new primes all split in $K$.  If
$s_2(E/K)=1$ and the nonzero addresses have product analytic order one and
uniformly bounded
$2j_h(k)-4w(h)-s(h)-s(hk)$, the base product has analytic order one.
\item The even coefficient system, exact same-filter square comparison,
divided operators, and trace identities of
Proposition~\ref{co:even} and Lemmas~\ref{co:operators},
\ref{co:traces}, and \ref{co:inertia} are available.
After a uniform lower bound for the relevant point indices has been proved,
the moving odd coefficient system and its stabilized minimum are available
as in Proposition~\ref{co:odd} and Lemma~\ref{co:minimum}.  The resulting unit symbols
satisfy the contraction, good-twin, and primary-isolation rules of
Proposition~\ref{co:symbol-interface}.  An even unit gives
\begin{equation}\label{pw:even-unit-interface}
 L_*(h)\ne0,\qquad v(L_*(h))\le 2w(h)+C;
\end{equation}
an odd unit gives analytic order one and
\begin{equation}\label{pw:odd-unit-interface}
 j(P_h)\le w(h)+C.
\end{equation}
The constants are independent of the number and sizes of the variable
primes.  The fixed genus partner used to define $P_h$ is denoted $h_*$.
\item Arithmetic realization and the forest and address statements of
Lemmas~\ref{gr:realization} and \ref{gr:forest} and
Theorem~\ref{gr:address} hold.
In particular, separate nonzero terminal polynomials may be addressed
simultaneously.  Primary auxiliaries preserve all fixed local classes,
have neutral incidence to specified stationary external vertices, and may
be required to split in any compatible fixed fields.  They contribute no
prime at address zero.
\end{enumerate}
\end{proposition}

\begin{proof}
These are the cited results, with the arithmetic ring-class hypotheses
verified by Proposition~\ref{rd:assembly}.  The even coefficient minimum
uses the cyclotomic lower bound; the odd detector of
Proposition~\ref{co:odd} still needs the arithmetic lower bound for
Lemma~\ref{co:minimum}.  Neither interpolation statement nor the even
symbol uses that odd normalization.
\end{proof}

The analytic witnesses and the rough completion used below are supplied by
Propositions~\ref{nv:local-witness} and \ref{nv:rough}; the latter has an absolute
prime-count bound.  The bounds in Proposition~\ref{pw:interfaces}, unlike the
eventual sizes of the primes chosen by nonvanishing or Chebotarev, must be
uniform in the parameters indicated there.

\begin{lemma}[Isogeny orientation and coranks]\label{pw:isogeny-corank}
Recall that a rational isogeny preserves analytic rank, Mordell--Weil rank,
and $2^\infty$-Selmer corank.  If $k$ is a quadratic discriminant and
$K=\mathbb Q(\sqrt{k})$, then
\begin{equation}\label{pw:quadratic-corank}
 s_2(E/K)=s_2(E/\mathbb Q)+s_2(E^{(k)}/\mathbb Q).
\end{equation}
For every $E/\mathbb Q$ one has, separately,
\begin{equation}\label{pw:kummer-corank}
 s_2(E)=r(E)+
 \operatorname{corank}_{\mathbb Z_2}\Sha(E/\mathbb Q)[2^\infty].
\end{equation}
\end{lemma}

\begin{proof}
These are Lemmas~\ref{pr:invariance} and \ref{pr:ranks}.  In particular,
the quadratic identity is an identity of coranks, not an integral
direct-sum decomposition.
\end{proof}

If the isogeny class contains a curve with full rational two-torsion, we use
that curve.  Otherwise choose a rational two-isogeny $E\to E'$ and write
$F=\mathbb Q(E[2])$, $F'=\mathbb Q(E'[2])$; both fields are quadratic.
A prime split in $F$ is called split, and one inert in $F$ is called simple.
Their weights are respectively $1$ and $1/2$.  A simple prime is good if it
is also inert in $F'$.  At such a prime
\begin{equation}\label{pw:good-simple}
 \#E(\mathbb F_p)\equiv2\pmod4.
\end{equation}
Indeed there is exactly one nonzero rational two-torsion point in the
reduction.  If the group order were divisible by four, its two-primary
subgroup would contain a rational point of order four.  Its image under the
two-isogeny would supply on $E'$ a rational two-torsion point different from
the dual kernel, contradicting inertness in $F'$.  Thus the good-twin rule
applies.  When $F=F'$, every simple prime is good.
A prime split in $F$ and inert in $F'$ is a \emph{donor}: attaching its
label to a bad simple label makes that label good.  This terminology
describes finite label operations; the required prime supplies will be
constructed below.

\subsection{Seeds and finite supplies}\label{pw:seeds-section}

The word ``label'' below means an individually realizable fixed-support
elementary Frobenius label, with orientations when required.  It does not
include a self-residue of a prime generator.  All contractions and
comparisons are those of Proposition~\ref{co:symbol-interface} and
Lemma~\ref{gr:forest}.

\begin{lemma}[Minimum seeds]\label{pw:minimum-seeds}
\begin{enumerate}
\item In an all-split alphabet, suppose a unit symbol has a nonempty
witness.  At the least positive vertex count $r$ of such a witness, all
mutual edge bits are immaterial.  The symbol is determined by the individual
fixed-support labels.  This remains true if the empty configuration is a
unit, provided no contraction to the empty configuration is used.
\item In the quadratic case, suppose the symbol has a nonzero witness,
the alphabet is closed under the allowed contractions, and it contains
the good twin operations.  There are
$r\ge0$ and $s_1\ge1$ such that for every $s\ge s_1$ of the permitted simple
parity, the minimum split count of a unit with $s$ simple vertices is $r$.
At these counts all edge bits are immaterial.  A chosen minimum seed may be
padded by arbitrarily many good twin pairs.
\end{enumerate}
\end{lemma}

\begin{proof}
For the first assertion, a derivative in an edge of an $r$-vertex graph is a
symbol on $r-1$ vertices.  If $r\ge2$, this graph remains nonempty and has
value zero by minimality.  If $r=1$, there is no edge.  The omitted-vertex
description then leaves exactly the fixed-support labels.

For the second assertion let $r(s)$ be the minimum split count at simple
count $s$, with value $+\infty$ when no witness exists.  Once a finite value
occurs, good-twin insertion gives $r(s+2)\le r(s)$.  A nonincreasing sequence
of nonnegative integers stabilizes.  Choose $s_1\ge1$ in that range.  A
derivative of an edge with a split endpoint reduces the split count and
preserves the simple count; it is therefore zero at the minimum.  A
simple--simple derivative is zero by the symbol rule.  Twin insertion gives
the last assertion.
\end{proof}

Call a list partitionable if some permitted contraction forest has a unit
symbol as its output.  A contraction certificate proves that its terminal
polynomial is nonzero: the indicated iterated finite difference is nonzero.

\begin{lemma}[Finite label reserves]\label{pw:reserves}
Fix a minimum seed and a finite label alphabet.
\begin{enumerate}
\item In the all-split rational alphabet, a list with the same total label
as the seed and sufficiently many copies of its labels is partitionable.
\item In the quadratic alphabet, retain $r$ split anchors with the seed
labels, and suppose the list and seed have the same total rational label
on the projection used by the symbol.  Suppose there are sufficiently
many simple vertices of the required parity and sufficiently many split
reserves of each realizable rational shift.  If $F\ne F'$, also suppose
\begin{equation}\label{pw:donor-balance}
 \Delta=\#\{p:\ p\text{ split in }F\text{ and inert in }F'\}
       -\#\{p:\ p\text{ inert in }F\text{ and split in }F'\}>C_0.
\end{equation}
Here $C_0$ depends only on the anchors and the reserves.  Then all other
split vertices can be attached to simple vertices so that the result is
the minimum seed and good twin pairs.  The twins may match a prescribed
finite enrichment of the rational labels.  The base simple labels may be
specified on any projection on which the required total agrees.
\end{enumerate}
All reserve requirements depend only on the finite alphabets, seed, and
projection.  They do not depend on the original list length.
\end{lemma}

\begin{proof}
In the first case retain the first $r-1$ seed labels and merge all other
vertices to the last cluster.  Sufficient copies ensure that this cluster
is nonempty.  Its label is forced by the total, hence is the last seed
label.  Lemma~\ref{pw:minimum-seeds} makes its unknown edges immaterial.

For the second case, the rational labels of split attachments form the
subgroup with trivial $F$-coordinate.  They act transitively on the simple
coset, so every prescribed change of a simple label has a realizable split
shift.  Reserve sufficiently many representatives of every such shift.
Remove these reserves and the anchors from the donor pool.  If $F\ne F'$,
\eqref{pw:donor-balance}, with $C_0$ exceeding the number removed, leaves a
distinct donor for each bad simple vertex.  Attach these donors; every
simple vertex is now good.  If $F=F'$, this operation was unnecessary.

Pair equal enriched good labels.  By withholding a bounded number of pairs,
leave between $s_1$ and $s_1+C_1$ unpaired vertices, with the correct parity;
$C_1$ depends only on the number of enriched labels.  Use reserved shifts
to turn the excess unpaired vertices into good pairs and to set the first
$s_1-1$ base positions to the desired seed labels.  The seed labels
themselves need not all be good.  Attach every remaining nonanchor split
vertex to the last base position.  Its projected label is forced by the
total, the anchors, the other base positions, and the twin pairs.  It is
the required last seed label.  Each cluster contains one simple vertex,
so every contraction is permitted.  Only a bounded number of reserve
shifts was used.

When incidences to a retained split anchor are needed only to remove twins,
they can be prescribed after the contracted labels are fixed, by the
independent terminal-edge realization.  Fixed stationary vertices instead
enter the finite label enrichment from the outset.  These two uses do not
increase the number of label types with the original list length.
\end{proof}

\begin{lemma}[Neutral packets]\label{pw:packets}
For any of the finite reserve problems above with prescribed rational total
zero, one can choose a finite packet of terminal labels with rational total
zero such that every nonempty union of copies of the packet has the required
supplies.  In the quadratic case with \(F\ne F'\), its donor surplus
may be made arbitrarily large.  When \(F=F'\), every simple label is
good and no donor surplus is required.  Compatible fixed splitting
conditions may be included whenever
the needed individual labels and shifts remain realizable.
\end{lemma}

\begin{proof}
Choose ample copies of every required label, including anchors and shifts.
In the quadratic case with \(F\ne F'\), add enough donors to exceed
all the losses from the
other labels and the prescribed reserve margin.  Duplicate every full
label in the resulting list.  Duplication gives total zero in every
rational elementary coordinate and preserves any strict donor surplus.
Choose the margins so that after reserving the finitely many anchors in
any one packet, the union of any positive number of packets still meets
the reserve requirements.  All labels are realized by distinct fresh
primes; duplication never means repeating an actual prime factor.
\end{proof}

\subsection{The even estimate and the corank-zero case}
\label{pw:lower-section}

\begin{proposition}[Partitionable even estimate]\label{pw:partitionable-even}
Fix an even symbol furnished by Proposition~\ref{pw:interfaces}.  For every
partitionable parameter $h$ on its filter with $s_2(E^{(h)})=0$,
\begin{equation}\label{pw:partitionable-bound}
 L_*(h)\ne0,\qquad v(L_*(h))-2w(h)-s(h)\le C.
\end{equation}
The constant depends only on the symbol and fixed local data.
\end{proposition}

\begin{proof}
If the parameter has no prime factors, a contraction certificate means
that the empty symbol itself is a unit.  Equation~\eqref{pw:even-unit-interface}
then proves the assertion directly.  Otherwise the stationary terminal
set is nonempty.
Keep the prime factors of $h$ stationary.  Their terminal polynomial is
nonzero by its contraction certificate.  The stationary one-symbol case
of Theorem~\ref{gr:address} gives arbitrarily large binary cubes based at
$h$, with only primary new primes, whose other addresses have unit symbol.
Equation~\eqref{pw:even-unit-interface} and $s(h_x)\ge0$ give the uniform
upper bound required by the cyclotomic transfer.  The base corank is zero,
so that transfer gives \eqref{pw:partitionable-bound}.  A base-dependent
dimension threshold does not affect the asserted constant.
\end{proof}

The original corank-zero base can now be treated without any odd
coefficient system or companion field.

\begin{proposition}[Even original base]\label{pw:even-base}
If $s_2(E)=0$, then $a(E)=0$.
\end{proposition}

\begin{proof}
Lemma~\ref{pr:invariance} gives even functional sign.  Fix an even symbol
on the filter of $1$, including all auxiliary class-group tests.  If its
value at $1$ is a unit, analytic nonvanishing is already detected.
Otherwise use a nonempty minimum seed in the full-torsion case, or a
stabilized quadratic seed in the other case.

For each $b$, assign a neutral packet from Lemma~\ref{pw:packets} to
every nonzero linear form on $\mathbb F_2^b$, and activate its primes
by that form.  At every nonzero address at least one packet is active;
their union satisfies the reserve problem and has a nonzero terminal
polynomial.  At zero no terminal is active.  The single-symbol case of
Theorem~\ref{gr:address} adds primary auxiliaries and makes the symbol a
unit at every nonzero address.  All local classes remain those of $1$.
Equation~\eqref{pw:even-unit-interface}, followed by cyclotomic transfer
at the original corank-zero base, gives $a(E)=0$.
\end{proof}

\subsection{The lower bound needed by the odd coefficients}
\label{pw:odd-lower-section}

For the odd construction, first fix the original coefficient and
class-group support and its unit tests.  Then choose the genus partner
$h_*$ as in Section~\ref{co:section}, adjoin its prime factors to the
support, and refine the unit tests there.  All variable $h$ avoid this
support.  The positive filter below is this prepared filter, fixed before
$h$ varies or the odd coefficients are normalized.

We next use the even estimate at analytically nonvanishing parameters,
not the corank-zero conclusion just proved.  The companion $k'=k'(h)$
in this step may vary with $h$, but its number of prime factors will be
bounded.  It is distinct from both the detector's fixed genus partner
$h_*$ and the fixed companion used in the final ring-class transfer.

At a parameter already known to have analytic rank zero, the estimate
\eqref{pw:partitionable-bound} transfers between either orientation of a
fixed two-isogeny.  The isogeny formula for the BSD quotient gives a bounded
change in $s(h)+2w(h)$.  At an active odd good prime the component-group
valuation is $t_p$, and the two active contributions are included in the
two weights.  At the fixed support only finitely many local twisting
classes occur.  Differential, period, and torsion indices are bounded by
the fixed isogeny and its reverse, and there is no regulator at rank zero.
This uses an isogeny identity for known finite groups, not BSD or converse
finiteness at an unknown parameter.

\begin{lemma}[A bounded companion for a varying odd parameter]
\label{pw:rough-companion}
Fix a positive odd-sign filter, including the coefficient and class-group
support used below, and a negative genus partner $h_*$ with
$L(E^{(h_*)},1)\ne0$, chosen using Proposition~\ref{nv:local-witness} and then
excluded from the variable parameters.  There are constants $A,C$ such
that every $h$ on this filter of analytic order one admits a negative odd
fundamental discriminant $k'$ for which
\begin{enumerate}
\item $(k',2Nh)=1$, $k'$ is a square at every place of $2Nh$, and
$\mathbb Q(\sqrt{k'})$ is nonexceptional;
\item $\omega(k')\le A$;
\item $L(E^{(hk')},1)\ne0$ and
\begin{equation}\label{pw:rough-bound}
 v(L_*(hk'))\le 2w(hk')+s(hk')+C.
\end{equation}
\end{enumerate}
\end{lemma}

\begin{proof}
At the fixed class-group support and the other fixed label tests, prepare
an even negative-index symbol for each possible total type at this
support, and for both isogeny orientations if needed.
There are finitely many such preparations.  The positive index $|hk'|$
has nonzero total rational type: its negative-discriminant parity is
already nontrivial among the tests at $2$.  Thus in the full-torsion case
a unit seed has positive count.  In the quadratic case use the stabilized
seed of Lemma~\ref{pw:minimum-seeds}.

When $F\ne F'$, interchanging the isogeny orientation interchanges the two
counts in \eqref{pw:donor-balance}.  Select the orientation with initial
balance at least zero.  The existing donors can then pay for arbitrarily
many existing bad simples.  They do not have to be supplied by new primes.

Let $R$ be the absolute prime-count bound in Proposition~\ref{nv:rough}.  Insert
in $k'$ a finite dummy supply that meets the anchors, minimum simple count,
and reserve requirements for every one of the finitely many prepared
symbols.  When \(F\ne F'\), add donor slack paying for every reserved
loss, the new dummy bad simples, and a further possible loss of $R$
from the later rough factor.  When \(F=F'\), all simple primes,
including the later rough primes, are already good; only the simple
and split-reserve supplies are needed.  Add a fresh ramified prime if
needed to exclude the exceptional
imaginary fields.  The number of dummy primes is bounded independently
of $h$.  Their labels and initial incidences may be chosen by
Lemma~\ref{gr:realization}.

Absorb $h$ and the dummy product into the curve.  Apply
Proposition~\ref{nv:rough} to complete $k'$ by a squarefree factor coprime to all
primes already used.  Prescribe its sign and its local unit classes so
that the final $k'$ is an odd negative fundamental discriminant, is a
square at $2Nh$, and has the required original local classes at $2N$.
These conditions give a nonempty reduced progression by the Chinese
remainder theorem.  Its root number is even: since $E^{(h)}$ has odd
sign and $k'$ is negative and square at every prime of $2Nh$,
the quadratic twist formula multiplies that sign by
$\chi_{k'}(-N_{E^{(h)}})=-1$.  Thus this is the positive-root-number
case of Proposition~\ref{nv:rough} for the absorbed curve.
Additional fixed primes may be excluded.  The proposition makes the
central value at $hk'$ nonzero with at most $R$ new factors.  Its size
threshold may depend on $h$ and all these congruences; $R$ does not.

No particular resulting unit class at the auxiliary class-group support
was assumed in advance: dummy supplies were prepared for every possible
total there.  For the actual total, Lemma~\ref{pw:reserves} makes $hk'$
partitionable.  Analytic nonvanishing gives Selmer corank zero by the
forward theorem.  Proposition~\ref{pw:partitionable-even} therefore gives
\eqref{pw:rough-bound}, with the isogeny comparison above if the other
orientation was used.  The finite preparation and the bound on dummy and
rough factors give $A,C$ independent of $h$.
\end{proof}

\begin{lemma}[Height comparison]\label{pw:height-comparison}
Here $j_h(k)$ is the index of the conductor-$|h|$ ring-class sum over
$\mathbb Q(\sqrt{k})$, whereas $j(P_h)$ is the index of the conductor-one
genus sum over $\mathbb Q(\sqrt{hh_*})$.
With the fixed modular parametrization and the fixed genus partner $h_*$,
the explicit Gross--Zagier formula gives, whenever the relevant factors
have analytic orders one and zero,
\begin{equation}\label{pw:height-ratio}
 2j_h(k)-2j(P_h)=v(L_*(hk))+O(1).
\end{equation}
For $k$ with bounded prime count as in
Lemma~\ref{pw:rough-companion}, the error is uniform in $h,k$; for a fixed
$k$ it is uniform in $h$.
\end{lemma}

\begin{proof}
Use the character-sum normalization of the explicit Gross--Zagier formula
in \cite[Theorem~1.1]{CaiShuTian2014}.  The restriction of $\chi_h$ to
$\mathbb Q(\sqrt{k})$ has exact ring-class conductor $|h|$: each odd
prime dividing $h$ splits and retains its ramified quadratic character.
The genus character over $\mathbb Q(\sqrt{hh_*})$ is unramified, so the
sum defining $P_h$ has conductor one.  Their canonical
heights are fixed modular factors times, respectively,
\[
 |h|\sqrt{|k|}\,L'(E^{(h)},1)L(E^{(hk)},1),
 \qquad
 \sqrt{|hh_*|}\,L'(E^{(h)},1)L(E^{(h_*)},1).
\]
In that theorem the multiplier is $c\sqrt{|D|}$, the level-discriminant
factor is $2^{\mu(N,D)}$, and the unit factor is
$[\mathcal O_c^\times:\mathbb Z^\times]^2$.  Here every level prime
splits, so $(N,D)=1$ and $\mu(N,D)=0$; the excluded extra-unit fields
make the unit factor $1$ in both applications.
Take their ratio.  The derivative and the rank-one regulator cancel;
the remaining varying real factor is $\sqrt{|hk|}$, precisely the factor
in $L_*(hk)$.  Comparison of the rational free line with its two quadratic
projections has bounded index because their composites are multiplication
by $2$.  The modular parametrization and the factor involving $h_*$ are
fixed.  The level is split and the discriminant and conductor are prime
to it; no growing conductor-prime product occurs in this normalization.
Local differential and period comparisons have finitely many possibilities
at the fixed support.  The extra-unit fields have been excluded, and the
number of additional places in the varying $k$ case is bounded.  These
observations give the asserted bounded valuation error.
\end{proof}

\begin{proposition}[Uniform lower bound for the odd detector]
\label{pw:odd-lower}
Every analytically rank-one $h$ in the prepared positive filter of
Lemma~\ref{pw:rough-companion} satisfies
\begin{equation}\label{pw:odd-lower-bound}
 2j(P_h)\ge 2w(h)+s(h)-C.
\end{equation}
In particular $j(P_h)\ge w(h)-C$ uniformly in $h$ and its prime count.
\end{proposition}

\begin{proof}
Choose $k'$ by Lemma~\ref{pw:rough-companion}.  Combine
\eqref{pw:ring-lower-interface}, \eqref{pw:height-ratio}, and
\eqref{pw:rough-bound}.  Since the factors are coprime,
$w(hk')=w(h)+w(k')$, and $w(k')\le\omega(k')\le A$.  Hence
\begin{align*}
 2j(P_h)
 &\ge 4w(h)+s(h)+s(hk')-v(L_*(hk'))-C\\
 &\ge 2w(h)+s(h)-2w(k')-C\\
 &\ge 2w(h)+s(h)-C'.
\end{align*}
Only even symbols and independently nonvanishing companions were used.
Thus this proof precedes, and does not assume, the odd minimum-depth
normalization.
\end{proof}

\subsection{From the moving odd coefficients to finite graphs}
\label{pw:odd-stages}

Proposition~\ref{pw:odd-lower} now supplies the arithmetic premise of
Lemma~\ref{co:minimum}.  We obtain an odd symbol with fixed minimum
normalized depth and the uniform unit bound
\eqref{pw:odd-unit-interface}.  To use the finite seed and packet lemmas,
we also need to stabilize its graph labels and values, not just this
numerical depth.

Keep the ultrafilter used in Lemma~\ref{co:minimum}.  The intrinsic fixed
support includes the genus partner $h_*$; at an actual stage it also
includes the detecting prime $\ell$.  That prime adds at most two places
in any fixed quadratic field.  The unit and class-group description of
$V_S(F)$ therefore bounds the dimension of the elementary label space
independently of the stage.  At $\ell$ the Schwartz function in
Lemma~\ref{co:theta-bridge} uses only the quadratic residue-line character:
increasing the Kummer precision does not impose a character of growing
order on the variable primes.  The old-support characters are fixed.

Choose bases and restrict to an ultrafilter-large set of stages on which
the finite label spaces, their addition tables, conjugation and projection
maps have one abstract form.  At every fixed weight bound there are only
finitely many typed graphs and symbol truth tables.  Select the table
occurring on an ultrafilter-large set.  These selected tables are
compatible under restriction and define one limit symbol.  Its contraction
and twin identities follow from the finite-stage identities.  The witness
sequence of Lemma~\ref{co:minimum} has bounded weight, hence takes only
finitely many graph types in this alphabet; one type is a witness in the
selected limit table.

If the empty configuration is a unit, a sufficiently advanced actual stage
already detects analytic order one at $h=1$.  Otherwise the limit witness
is nonempty, so Lemma~\ref{pw:minimum-seeds} supplies the minimum seeds
and stabilized simple count needed below.  Fix them in this remaining
case; Lemmas~\ref{pw:reserves} and \ref{pw:packets} apply to their finite
label alphabets.  Every
particular contraction or twin comparison uses only a bounded graph.
Its selected truth table agrees with the actual symbol on an
ultrafilter-large set of stages.  Thus any prescribed finite collection
of such comparisons holds at one sufficiently advanced stage.

Later we will add the fixed support of a separate companion and the label
projection read by its even test.  Once compatible lifts have been
verified, we may stabilize this finite enrichment by restricting the
\emph{same} ultrafilter.  This preserves the intrinsic odd symbol, its
minimum, and its chosen seeds.  The actual stage is selected only after
the cube dimension and all its finite networks have been chosen: their
maximum weight gives a finite list of table and precision requirements.
No intersection over all weight bounds, and no single stage working for
every cube dimension, is required.

\subsection{A nonempty split-field symbol with full two-torsion}
\label{pw:split-section}

We give the extra argument required when $E[2]$ is rational.  It addresses
the possibility that the usual even minimum is attained only at the empty
configuration after restricting all new primes to split in the companion
field.

\begin{lemma}[Recurrence congruences]\label{pw:recurrence}
Let $p$ be odd, $a\equiv p+1\pmod4$, and let
\[
 b_0=1,\quad b_1=a-\kappa,\quad
 b_{s+1}=ab_s-pb_{s-1}.
\]
If $\kappa=\pm1$, then $b_{s+1}-b_s$ is even and
\[
 \frac{b_{s+1}-b_s}{2}\equiv\frac{a-\kappa-1}{2}\pmod2
 \quad(s\ge0).
\]
If $\kappa=0$, put $b_{-1}=0$; then $b_s-b_{s-1}$ is odd for every
$s\ge0$.  Finally, for $\kappa=1$, if $a\equiv p+1\pmod{2^m}$ then
$b_s\equiv p^s\pmod{2^m}$ for every $s\ge0$.
\end{lemma}

\begin{proof}
Modulo $2$ the recurrence is $b_{s+1}=b_{s-1}$.  Its initial values prove
the parity assertions.  For $d_s=b_{s+1}-b_s$, the recurrence for $d_s$
is the same.  Directly,
$d_1-d_0\equiv(p-1)(p-\kappa)\equiv0\pmod4$ when $\kappa=\pm1$.
Thus $d_s/2$ is constant modulo $2$.  The last assertion follows by
substituting the sequence $p^s$ in the recurrence and checking its first
two terms.
\end{proof}

\begin{lemma}[Nonempty $K$-split witness]\label{pw:split-witness}
Suppose $E[2]$ is rational.  Choose a negative unit $k$ of a fixed even
symbol for $E$, prime to all previous fixed support and to the ramification
support of $E[8]$.  Put $K=\mathbb Q(\sqrt{k})$, chosen nonexceptional.
Let $S_0$ contain the complete fixed support of the coefficient system,
$2N$, the primes of $k$, and primes generating the class group of $K$.
On the positive filter of $1$ square at $S_0$, there is an
even unit test on nonempty products of primes split in $K$ with these
properties:
\begin{enumerate}
\item it has a nonempty witness;
\item it satisfies the oriented split contraction and primary-isolation
rules over $K$, for contractions with nonempty output;
\item a unit at $D$ implies
\begin{equation}\label{pw:split-unit-bound}
 L(E^{(kD)},1)\ne0,\qquad v(L_*(kD))\le 2\omega(D)+C_k.
\end{equation}
\end{enumerate}
This conclusion uses the integral trace and local identities in
Proposition~\ref{pw:interfaces}(3).
\end{lemma}

\begin{proof}
If the ordinary coefficient test already has a nonempty split witness,
there is nothing to construct.  Otherwise we first obtain an additional
trace congruence modulo $8$.  That congruence makes one further division
of the coefficient form integral.  We verify its contraction identities,
then use a prime inert in $K$ to force a nonzero test at a split prime.

\emph{Extra depth if the ordinary test has no witness.}
Use the ordinary even coefficient form for $E^{(k)}$, normalized by
$a(1)=1$.  The positive filter is chosen so that $kD$ lies in the original
negative filter.  Since $k$ attains its least normalized depth, the exact
ratio comparison of Proposition~\ref{co:even} gives
\begin{equation}\label{pw:split-basic-depth}
 v(a(D))\ge\omega(D)
\end{equation}
at every allowed squarefree $D$.  The factor $|k|$ cancels from this ratio;
one is not subtracting unrelated bounded-error inequalities.

If equality holds at a nonempty product of $K$-split primes, use the usual
even bit, projected from oriented $K$-labels to rational labels.  It has
all the required properties by Proposition~\ref{co:symbol-interface}.  We
therefore assume for the rest of the proof that equality never occurs
there.  By the rational-phase assertion in Proposition~\ref{co:even},
normalization by $a(1)=1$ makes each nonzero squarefree coefficient a
rational number times a root of unity.  Its valuation is an integer,
not just an element of the possibly ramified coefficient field's
discrete value group.  Thus strict inequality implies
\begin{equation}\label{pw:split-extra-depth}
 v(a(D))\ge\omega(D)+1
 \quad\text{for every nonempty allowed all-$K$-split }D.
\end{equation}

Write $g=\sum a(n)q^n$ for the form.  Its characters satisfy
$\psi/\eta=\chi_{-4}$, so $\eta$ is odd.  For a set $U$ of $K$-split
primes put
\[
 B_U=2^{-|U|}\prod_{p\in U}(U_p-\eta(p)V_p)g,
 \qquad
 \Theta_J=\sum_{m>0,\ (m,S_0)=1}\eta(m)m q^{Jm^2}.
\]
All expansions retain the prescribed support sieves.  The ordinary divided
form is
\[
 H_U=\frac12\left(B_U-\sum_{J\subseteq U}B_U[J]\Theta_J\right).
\]
Here $B_U[J]$ denotes its coefficient at the squarefree index $J$.

\emph{The resulting trace congruence.}
We first prove
\begin{equation}\label{pw:split-mod8}
 a_p(E^{(k)})\equiv p+1\pmod8
 \quad\text{for every allowed good }p\text{ split in }K.
\end{equation}
Apply the ordinary $H_{\{p\}}$ trace at a Frobenius lift $\phi_p$,
testing coefficient $p^2$ after theta multiplication.  The form
$H_{\{p\}}$ is already integral, and $\theta\equiv1\pmod2$;
hence theta multiplication does not alter any of its reduced
coefficients after these divisions.  On the
conductor-one ramified principal-series part the identity is
\[
 U_p^2t(\phi_p)=U_p^3+d(\phi_p)U_p.
\]
On the unraised unary part it is the ordinary good-prime coefficient
identity.  These are the only two types in this expression.  Consequently
\[
 (t(\phi_p)f)[p^2]=f[p^3]+(d(\phi_p)f)[p].
\]
Replace the trace by a fresh split prime $p'$ with the same finite data,
including every support-unit test.  At each rational quadratic character
unramified outside $S_0$, a Frobenius lift at $p$ has the character value
of $p$.  Hence the matched $p'$ has the same value, and $pp'$ has value
$1$.  Equivalently, $p'$ has the same unit squareclasses as $p$ at all
places of $S_0$, so their product is square at all of them, including $2$.
Since $\phi_p\in G_K$, the new prime is also split in $K$.
Its coefficient on the left tests $pp'$; by
\eqref{pw:split-extra-depth} it is zero modulo the maximal ideal after
the two divisions.  The determinant term on the right is zero: the
unary summand at $p$ has the same character and cancels there.

On the right, the indices $p$ and $p^3$ in $B_{\{p\}}$ have underlying
squarefree index $1$; they involve only $a(1),a(p^2),a(p^4)$.  The
$J=\varnothing$ unary term has no coefficient at either nonsquare
index.  Thus no other squarefree coefficient contributes there.
For $a=a_p(E^{(k)})$, the phase-free square coefficients are
$1,a-1,a(a-1)-p$.  Hence the remaining coefficient is, up to a unit,
\[
 H_{\{p\}}[p^3]
 =\frac{a^2-2a-p+1-p(a-2)}4
 =\frac{(a-1)(a-p-1)}4.
\]
It is zero modulo $2$.  Full rational two-torsion makes $a$ even, so
$a-1$ is odd; this proves \eqref{pw:split-mod8}.  No individual total
filter condition on $p$ was needed: its matched prime $p'$ makes $pp'$
lie on the filter, and the square-index terms on the right come from $1$.

\emph{One further integral division.}
Now define
\begin{equation}\label{pw:four-divided-form}
 \widetilde H_U=\frac14
 \left(B_U-\sum_{J\subseteq U}B_U[J]\Theta_J\right).
\end{equation}
We verify integrality at every coefficient.  Write a test index uniquely
as $Jrm^2$, where $J\subseteq U$ and $r$ is squarefree and disjoint from
$U$.  The underlying squarefree coefficient is at
$D=(U/J)r$.  Removing the unit phases, a prime of $U/J$ contributes an odd
factor $b_s-b_{s-1}$ with $\kappa=0$, and a prime of $J$ contributes
$b_{s+1}-b_s$ with $\kappa=(D/p)$.  Thus
\begin{equation}\label{pw:shell-valuation}
 v(B_U[Jrm^2])\ge\omega(r)
\end{equation}
by \eqref{pw:split-basic-depth}, with one additional unit of valuation
when $D$ is nonempty and all its primes split in $K$.

If $r\ne1$ and all primes of $D$ split in $K$, the last inequality is
at least $\omega(r)+1\ge2$.  If $D$ contains an inert prime, then $r$
contains at least two: the total rational $K$-character is trivial on
the filter, and all primes of $U$ split.  Then
\eqref{pw:shell-valuation} is again at least $2$.  No unary term has
this squarefree part.  This proves divisibility by $4$ for $r\ne1$.

Suppose $r=1$ and $J\ne U$.  Then $D=U/J$ is nonempty and all split,
so $B_U[J]$ is even.  This cancellation can be checked without dividing
by a leading multiplier, which may be zero.  Factor from both shell
coefficients the common phase and the scalar
$2^{-|U|+|J|}a(U/J)$, whose valuation is at least $1$ by
\eqref{pw:split-extra-depth}.  The remaining factors are the odd
differences at $U/J$, the integral divided differences at $J$, and
the ordinary square-index factors outside $U$.  By
Lemma~\ref{pw:recurrence}, their product is congruent modulo $2$ to
the product of their leading values, even when one of those values
is zero modulo $2$.  Since $m$ is odd, multiplication by $m$ preserves
this congruence.  The common scalar supplies the second factor $2$.
Retaining the phase $\eta(m)$ therefore gives
\[
 B_U[Jm^2]\equiv B_U[J]\eta(m)m\pmod4.
\]
This is exactly the cancellation with $B_U[J]\Theta_J$.

Finally let $J=U$, so $D=1$.  At a prime outside $U$ the recurrence gives
$b_s\equiv p^s\pmod4$.  At a prime of $U$,
\eqref{pw:split-mod8} gives $b_s\equiv p^s\pmod8$, and hence
\[
 \frac{b_{s+1}-b_s}{2}
 \equiv \frac{a_p-2}{2}\,p^s\pmod4.
\]
Multiplying these congruences, again retaining the phases, gives
$B_U[Um^2]\equiv B_U[U]\eta(m)m\pmod4$.  This proves integrality in
the last case and establishes \eqref{pw:four-divided-form}.

For the trace argument we specify the integral module being used.
Write $\psi_U$ for the character of $B_U$.  The character of $\Theta_U$
is also $\psi_U$, so its diamond correction is zero.  For the proper
subsets the diamond formula is
\[
 \langle u\rangle\widetilde H_U
 =\psi_U(u)\widetilde H_U+
 \sum_{J\subsetneq U}\frac{B_U[J]}2
 \frac{\psi_U(u)-\eta(u)(-J/u)}2\Theta_J.
\]
Each $B_U[J]/2$ here is integral, and each divided character difference
is zero or a unit.  Let $\widetilde{\mathcal M}_U$ be the $\mathcal O$-module
generated by $\theta\widetilde H_U$ and
$\theta(B_U[J]/2)\Theta_J$ for $J\subsetneq U$, together with all their
good-Hecke and diamond translates.  Its generators are integral by this
formula and the good-Hecke coefficient formula.  The finite-module and
Chebotarev-continuity argument of Lemma~\ref{co:traces} therefore applies
to $\widetilde{\mathcal M}_U$.  A fresh trace kills every displayed unary
correction after coefficientwise reduction, also after further fresh
tests.  Thus the same coefficientwise reduced determinant and trace
identities hold on the required orbit.

\emph{The contraction and isolation identities.}
On $\theta\widetilde H_U$, two fresh traces in $G_K$ give zero after
reduction at every coefficient.  To see this, start with an arbitrary
test index $Jrm^2$ and choose the two fresh primes outside that index
and the level.  The filter imposes $\chi_K(D)=1$; since $U$ is all split,
the number of inert factors in $r$ is either zero or at least two.
The fresh traces replace $r$ by $rq_1q_2$, with both $q_i$ split.
In the first case the valuation after division by $4$ is at least
$\omega(r)+2+1-2=\omega(r)+1\ge1$.  In the second case it is at least
$\omega(r)+2-2=\omega(r)\ge2$.  The unary terms disappear under these
fresh tests.  Additional fresh $G_K$ traces only improve these estimates,
so vanishing holds on the entire reduced trace orbit.  The trace identity
therefore gives
\[
 t(gu^2)=t(g)\qquad(g,u\in G_K)
\]
on this reduced trace orbit.  Thus the tests factor through the elementary
extension of $K$ unramified off the level.

A single fresh $K$-split trace at coefficient $1$ reads the bit
\begin{equation}\label{pw:deeper-bit}
 \widetilde A(D)=
 \left(2^{-\omega(D)-1}a(D)\text{ with its phase removed}\right)\bmod2
\end{equation}
for a nonempty all-split configuration.  The contraction rule follows
from the same integral local identity as in Lemma~\ref{co:inertia}.
For completeness its first shell, with $p\in U$ and a fresh trace $q$,
has underlying configuration $D'=(U/p)q$, which is still nonempty, and
its reduction is
\[
 c_p(D')\widetilde A(D'),\qquad
 c_p(D')=\frac{a_p-(D'/p)-1}{2}\bmod2.
\]
Changing the oriented inertia bit changes $c_p$ once.  The omitted-prime
identity gives the product label and removes the entire rational prime
$p$ from the new ramification support.  Consequently its repeated
contraction has precisely the forest interpretation of
Lemma~\ref{gr:forest}.  An omitted-vertex label equal to the identity has
trace zero, which is primary isolation.  No contraction to the empty
configuration has occurred.

\emph{A nonempty witness for the deeper test.}
It remains to prove that this deeper test has a witness.  Fresh
ramification gives $K\not\subseteq\mathbb Q(E[8])$.  By Chebotarev there
is a good prime $p$ inert in $K$ and trivial on the original $E[8]$.
Then
\[
 p\equiv1\pmod8,\qquad a_p(E^{(k)})\equiv-2\pmod8.
\]
On $\widetilde H_\varnothing$, whose coefficient at $1$ is zero, test
$t(u)^2$ at coefficient $1$, where $u=\operatorname{Fr}_p$.  The square
coefficient and matching unary character give, up to a unit,
\[
 \frac{a_p(E^{(k)})-1-p}{4}\bmod2=1.
\]
The determinant term at coefficient $1$ cancels.  In characteristic two,
$t(u^2)=t(u)^2$, so a single trace in $G_K$ is nonzero.  A fresh prime
realizing this finite trace test is split in $K$.  Moreover $u^2$ is
trivial on every rational quadratic character, so the prime lies on the
required total rational filter.  It is a nonempty witness for
\eqref{pw:deeper-bit}.  Finally the exact coefficient ratio gives
\eqref{pw:split-unit-bound}, with the additional depth at most one.
\end{proof}

\begin{lemma}[Orientation correction]\label{pw:coinvariants}
Let $K$ be imaginary quadratic, let $S_0$ contain $2$ and its ramification
support, and let $V_0$ be the elementary label space over $K$ unramified
off $S_0$.  Complex conjugation acts on $V_0$ by $c$.  If the rational
total of an oriented prime configuration is zero on every rational
quadratic character unramified off $S_0$, then its oriented total belongs
to $(1+c)V_0$.  A reserved supply containing every label allows its
orientations to be flipped to realize any prescribed change in
$(1+c)V_0$.
\end{lemma}

\begin{proof}
Complex conjugation has order two outside $G_K$, so
$G_{\mathbb Q}=G_K\rtimes\langle c\rangle$.  A $c$-invariant quadratic
character of $G_K$ extends by giving $c$ value zero.  At a rational prime
outside $S_0$, inertia lies in $G_K$ and the character is trivial there;
thus its extension remains unramified outside $S_0$.  The assumed total
therefore annihilates $(V_0^*)^c$.

For a finite vector space with involution, the annihilator of
$\ker(1+c^*)=(V_0^*)^c$ is $\operatorname{im}(1+c)$.  This proves the
first assertion.  Flipping a vertex of label $v$ changes its oriented
label from $v$ to $cv$, and hence changes the total by $(1+c)v$.
Reserve one representative for every preimage needed to express the
desired element of the image.  Their independent flips give the second
assertion.  These flips preserve every rational norm label.
\end{proof}

\subsection{The corank-one original base}\label{pw:assembly-section}

\begin{proposition}[Odd original base]\label{pw:odd-base}
If $s_2(E)=1$, then $a(E)=1$.
\end{proposition}

\begin{proof}
Parity gives odd functional sign.  Use the normalized odd limit symbol
on the positive filter of $1$ from Section~\ref{pw:odd-stages}, with
its fixed support, genus partner $h_*$, and chosen ultrafilter.  If it
detects the base, a sufficiently advanced actual stage proves the result.
Otherwise choose its nonempty minimum seed.  The moving detecting prime
$\ell$ is not an extra fixed support requirement.

\emph{Fixing the companion.}
Choose a separate even symbol for negative discriminants square at $2N$,
excluding the fixed odd support, the relevant torsion ramification, and
the exceptional imaginary fields.  Choose a unit $k$ of that symbol.
Then $L(E^{(k)},1)\ne0$.  Fix this $k$ and
$K=\mathbb Q(\sqrt{k})$ for the rest of the proof, before choosing any
cube dimension.  Its fresh ramification makes it independent of the old
fixed label extensions.  At sufficiently advanced stages $\ell$ avoids
$k$.  The intrinsic odd label tests continue to omit the primes of $k$.

We construct common terminal packets, individually split in $K$, such
that the odd terminal polynomial and an even terminal polynomial testing
$hk$ are separately nonzero at each nonzero address.  They need not have
a common unit evaluation yet; the address lemma will supply that step.

\emph{The quadratic residual case.}
First suppose $F,F'$ are quadratic.  Keep the original even symbol
evaluated on $hk$, treating the vertices of $k$ as stationary external
data.  Enrich the rational labels of the $h$-vertices by the primes of
$k$.  Their total there and at the even support is required to be zero.
Let $V$ be the resulting rational elementary label space, including the
$F,F'$ and $K$ coordinates and all relations among these characters.
The allowed simple labels form the affine set
$\{v\in V:\chi_F(v)=1,\ \chi_K(v)=0\}$, in additive notation.
The difference of any two of its elements belongs to
$\ker\chi_F\cap\ker\chi_K$, precisely the realizable rational labels
of the permitted split attachments.  Translation by this subgroup is
therefore transitive.  The fields $F,F'$ are unramified outside the old
support, whereas $K$ has a fresh ramified prime; thus $K$ is independent
of their compositum.  When $F\ne F'$, this realizes both the good-simple
pattern $(1,1,0)$ and the donor pattern $(0,1,0)$ in their three
coordinates.  When $F=F'$, all simple labels are good.

For the old odd projection, take a normal closure of its finite label
extension at any stage with $\ell$ avoiding $k$.  It is unramified outside
the old support and $\ell$.  Once $\ell$ avoids
$k$, a prime ramified in $K$ remains outside that set.  Hence $K$
has trivial intersection with this normal closure, and every old odd
label can be realized together with splitting in $K$.  This applies
also to oriented split anchors over $F$.  The added stationary $k$
coordinates may have relations, but they are already included in $V$;
they do not impose extra requirements on the omitted odd projection.
These compatible lifts let us stabilize the common alphabet and its
projection maps on the same ultrafilter, as in Section~\ref{pw:odd-stages}.
The intrinsic odd minimum and seed are unchanged.  Choose the following
packets in this fixed abstract alphabet, before selecting an actual stage.

Use packets with sufficiently large simple supply and reserves in this
kernel, and with donor slack when \(F\ne F'\).  For the even test, take
no split anchors.  If the fields are distinct, correct all bad simples
to good ones; if they coincide, every simple is already good.  Pair
equal enriched labels, and
retain a nonempty good base pair.  Use reserves to match the pair and
the remaining good pairs, and attach every leftover split vertex to the
last base position.  Its label is forced by total zero.  The output is
a collection of good twin pairs, each matching also at every split
vertex of $k$.  Twin insertion into the unit at $k$ proves that this
output has even symbol one.  Thus the even terminal polynomial is
nonzero.

For the odd test, project away the newly added tests at $k$.  The same
packets contain enough split anchors, simple vertices and reserve shifts
to contract to the stabilized odd seed and good pairs, by
Lemma~\ref{pw:reserves}.  Its total on that projection is the prescribed
zero total.  Incidences between the retained anchors and the good pairs
can be set after contraction, by independent edge choices.  The odd
polynomial is therefore nonzero as well.  This odd test does not include
the stationary vertices of $k$.

\emph{Full rational two-torsion.}
Now suppose the curve has full rational two-torsion.  Use the even test
of Lemma~\ref{pw:split-witness}, with $K$-oriented labels on its support
$S_0$, and let $v_1,\ldots,v_r$ be a minimum nonempty witness.  Its
oriented total $v_1+\cdots+v_r$ lies in $(1+c)V_0$, by
Lemma~\ref{pw:coinvariants}.  This even test uses the oriented elementary
extension of $K$, whereas the odd test reads only its rational projection;
the even contraction rule is needed only for nonempty outputs.
Every odd seed label lifts to a $K$-split prime because $K$ was freshly
ramified outside its intrinsic support.  Each even label also has a
compatible lift to the common alphabet.  Stabilize the finite common
alphabet and projections on the same ultrafilter, preserving the
intrinsic odd symbol as in Section~\ref{pw:odd-stages}.

Choose packets containing ample copies of every needed even label, all
odd seed labels on their rational projection, and orientation-shifting
reserves.  Arbitrary pairs of projected labels need not lift: take a
union of sufficiently many compatible lifts for each projection
separately.  Duplicate each full label, so the packet has rational
total zero.

For the even contraction, retain anchors with labels
$v_1,\ldots,v_{r-1}$.  After reserving these anchors, use orientation
flips to make the total equal to $\sum v_i$; this is possible by
Lemma~\ref{pw:coinvariants}.  In a contraction certificate these flips
are choices of conjugations along the forest, as permitted by
Lemma~\ref{gr:realization}; they do not change the common terminal data
separately at different addresses.  Merge all remaining vertices into the
last cluster.  Its label is $v_r$.  The cluster is nonempty, and the
minimum nonempty count makes the output a unit independently of its
mutual edges.  For the odd contraction, use its rational projection.
Retain all but the last odd seed label and merge the remainder.  Total
zero gives the last label,
so this is also a unit.  Orientation corrections for the even test did
not change these rational labels.

\emph{Simultaneous addressing and an actual stage.}
Now fix $b$ and put a sufficient common packet at every nonzero linear
form on $\mathbb F_2^b$.  Each nonzero address has two separately
nonzero terminal polynomials.  In the fixed abstract alphabet,
Theorem~\ref{gr:address} gives finite primary-auxiliary networks making
both symbols one there.  Require their primary types to split in $K$
(and $F$ in the quadratic case) and to be neutral at the stationary
primes of $k$.  They are primary at every support coordinate used by
either projection, so both symbol identities apply to the same common
graph data.

Let $W_b$ bound all graph weights required by these finite networks.
A finite intersection of ultrafilter-large sets supplies an actual stage
with the same truth table through $W_b$, all
required divisibilities and trace identities, and sufficient precision
for the prescribed divisions.  The moving prime also avoids the fixed
companion $k$.  Realize the resulting finite graph at that stage using
Lemma~\ref{gr:realization}.  All active products $h$ are positive and
square at the fixed tests, and all their primes split in $K$.  No
auxiliary appears at address zero.  The depth constant and fixed
companion remain unchanged as $b$ grows.

\emph{Returning to the original curve.}
At every nonzero address the analytic orders of $E^{(h)}$ and $E^{(hk)}$
are respectively one and zero.  The odd unit bound, the even unit bound
(or \eqref{pw:split-unit-bound}), and
Lemma~\ref{pw:height-comparison} with fixed $k$ give
\[
 2j_h(k)\le 4w(h)+C.
\]
The finite Sha lengths at these addresses are nonnegative, so this also
bounds $2j_h(k)-4w(h)-s(h)-s(hk)$ from above.

At address zero, $L(E^{(k)},1)\ne0$ gives $s_2(E^{(k)})=0$ by the
forward theorem.  Lemma~\ref{pw:isogeny-corank} therefore gives
$s_2(E/K)=1$.  Ring-class missing-vertex transfer now makes the base
product $L(E,s)L(E^{(k)},s)$ have a simple zero.  Its second factor is
nonzero at $1$, and hence $a(E)=1$.
\end{proof}

\begin{proof}[Proof of Theorem~\ref{thm:main}]
Use Lemma~\ref{pw:isogeny-corank} to choose the indicated isogeny
orientation before constructing either coefficient system.
The even coefficient system and cyclotomic transfer give both
Proposition~\ref{pw:partitionable-even} and the corank-zero conclusion
of Proposition~\ref{pw:even-base}.  The bounded varying companion,
ring-class lower bound, and height comparison then establish
Proposition~\ref{pw:odd-lower}.  This discharges the arithmetic premise
before the odd normalization and finite-graph construction of
Section~\ref{pw:odd-stages}.

For the corank-one case, Proposition~\ref{pw:odd-base} fixes a nonvanishing
even companion before the cube dimension, constructs simultaneous unit
vertices, and applies ring-class transfer at the original base.  This gives
$a(E)=s_2(E)$ in both cases.  The established Gross--Zagier--Kolyvagin implication in
analytic rank zero or one gives $r(E)=a(E)$ and finiteness of the whole
Shafarevich--Tate group.  Equation~\eqref{pw:kummer-corank} remains the
underlying corank identity throughout.  In particular, neither a finite
two-Selmer dimension nor Mordell--Weil rank was substituted for the
original Selmer corank hypothesis.
\end{proof}

\begin{remark}[The prescribed base]\label{pw:status}
The final companion has Selmer corank zero because it was chosen with
nonzero central value.  Thus \eqref{pw:quadratic-corank} supplies the
required corank over $K$; it is not an additional companion hypothesis.
Only the finite realizations and their precisions change as the cube
grows.  The valuation constants remain fixed, so transfer applies to the
prescribed original curve without deleting a density-zero set of bases.
\end{remark}
